\section{The actual wave increment and the source's three compatibility targets}
\label{momop:section}

The source is the retained 166-page manuscript. Its SHA-256 is given
in these two consecutive segments:
\[
\begin{gathered}
\texttt{0e779481c4da40bd28d1e642e1d8ca57}\\
\texttt{447d129610df28dfa5a11e9af8ae228f}.
\end{gathered}
\]
The source operations used here are (8.3)--(8.27) and the complete
four-step cycle of Proposition 9.6, on printed pages 89--111.
Its coefficient classes are Definitions 6.4--6.5 and Proposition 6.6,
printed pages 69--72. The source estimates quoted below are attributed
to those statements. The following calculation of the actual candidate
targets and all pressure-induced changes is a direct finite calculation.
It does not infer source iteration exponents from finite smoothness.
The source correction cycle is its viscosity-one cycle. The physical
increment identities below separately retain any fixed
$\nu_{\rm NS}>0$; using another physical viscosity in those identities
does not identify the resulting candidate with the source's witness.

\subsection{The actual covariance increment before pressure reconstruction}

Use right-handed cylindrical coordinates $(r,\theta,z)$. An overbar
means the normalized angular average followed by normalized auxiliary
Haar average, taken at fixed slow variables before phase evaluation.
The old source state has actual zero-angular-mean wave $w_*$, including
all its curl corrections. The new candidate is the proved finite field
\[
 w_y=\sum_{j,a}y_a(t_{c,j})Z_{ja},\qquad
 \pi_y=\sum_{j,a}y_a(t_{c,j})\Pi_{ja},\qquad
 \langle w_y\rangle_\theta=\langle\pi_y\rangle_\theta=0.
\]
The $Z_{ja},\Pi_{ja}$ are the complete physical shapes from the
third-growth quantitative calculation; $y_a$ is the actual coupled
pair, with $t_{c,j}=t_b+(r_a/L_{s,j})v_j$. No amplitude factor is
removed from an auxiliary average: it can depend on $v_j(Y)$.
Write the exact double-averaged covariance increment as
\begin{equation}\label{momop:C}
 C_{ab}=
 \overline{w_{*,a}w_{y,b}+w_{y,a}w_{*,b}+w_{y,a}w_{y,b}},
 \qquad a,b\in\{r,\theta,z\}.
\end{equation}
Equivalently, substitute the displayed finite sum for both occurrences
of $w_y$ inside this average. In particular $C$ retains every old-wave
cross term and every same-label quadratic term. It is symmetric.
All its radial integrals below are finite because the new fields are
smooth and supported in a compact interior annulus. The old field
need only be restricted to this support.

The old source finite correction state has the two exact mean-velocity
constraints
\[
 \int_0^\infty r^2\overline v\,dr=0,\qquad
 \int_0^\infty r\overline\gamma\,dr=0 .
\]
These are source (9.10) in physical variables. The present initial
addition $w_y$ leaves $(\beta,v,\gamma)$ unchanged because its angular
mean is zero. The temporal mean operation has zero auxiliary mean,
and the first two rows of the five-row operation below give zero
increments of these integrals. Thus all specified operations preserve
their initially zero values. Preservation, by itself, does not change
a nonzero entering value to zero.

The already calculated raw tangential residual moments are
\[
 M_2=\partial_zJ_2,\qquad M_1=\partial_zJ_1,\qquad
 J_2=\int_0^\infty r^2C_{z\theta}\,dr,\quad
 J_1=\int_0^\infty rC_{zz}\,dr.
\]
Here the candidate has been added with its zero-angular-mean pressure
$\pi_y$, so the mean pressure is still the old mean pressure.
The source algorithm next reconstructs the mean pressure from the
updated velocity. This operation changes the axial residual moment.

To compute it, first let $\mathcal C_{ab}$ denote the angular average
of the tensor inside \eqref{momop:C}, retaining its auxiliary dependence.
Let $\mathcal D_r$ be the physical radial derivative on the extended
domain, including the auxiliary phase-map derivative. Its auxiliary
average is ordinary $\partial_r$. Since the mean velocity is unchanged
by $w_y$, subtracting the radial conservative equations gives
\begin{equation}\label{momop:delta_g}
 \Delta g_r=-(\mathcal D_r+r^{-1})\mathcal C_{rr}
                   -\partial_z\mathcal C_{zr}
                   +r^{-1}\mathcal C_{\theta\theta}.
\end{equation}
Indeed the radial transport of a divergence-free velocity is
$(\mathcal D_r+r^{-1})u_r^2+
r^{-1}\partial_\theta(u_\theta u_r)+
\partial_z(u_zu_r)-u_\theta^2/r$.
Angular averaging removes the angular derivative and all products
containing exactly one old mean factor and one zero-angular-mean wave.
The source defines $g_r$ as the negative of the mean radial residual
apart from the pressure gradient, which gives exactly the signs in
\eqref{momop:delta_g}. The time and viscous averages of $w_y$
are zero, so no physical viscosity has been dropped in this subtraction.
The formula remains true for any fixed $\nu_{\rm NS}>0$ in the
physical residual.

Set $g=\overline{\Delta g_r}$. Auxiliary derivative integrals vanish,
and hence
\begin{equation}\label{momop:gbar}
 g=-(\partial_r+r^{-1})C_{rr}-\partial_zC_{zr}
                                    +r^{-1}C_{\theta\theta}.
\end{equation}
The increment of the source's pressure defect is therefore the
following explicit new scalar, rather than either raw tangential moment:
\begin{equation}\label{momop:deltaP}
 \Delta P=\int_0^\infty g\,dr
 =\int_0^\infty\frac{C_{\theta\theta}-C_{rr}}r\,dr
                           -\partial_z\int_0^\infty C_{zr}\,dr .
\end{equation}
The equality integrates the derivative $\partial_rC_{rr}$;
both endpoints vanish by compact support. The weighted second
moment of this same source is
\begin{equation}\label{momop:gsecond}
 \int_0^\infty r^2g\,dr
 =\int_0^\infty r(C_{rr}+C_{\theta\theta})\,dr
                   -\partial_z\int_0^\infty r^2C_{zr}\,dr .
\end{equation}
For its first term, integration by parts gives
$-\int r^2\partial_rC_{rr}=2\int rC_{rr}$, after which the
explicit $-C_{rr}/r$ contribution subtracts one copy.
These two identities retain the radial and axial coupling exactly.

\subsection{Compact mean pressure and the corrected axial flux moment}

Retain the source concentration scale and its full axial derivative:
\[
 z=q^D\eta,\qquad 1-t=q(1-\eta^2),\qquad
 A=\tfrac12+h,\quad D=\tfrac12-h,\quad
 q_z=\frac{2\eta q^{1-D}}{1-2h\eta^2}.
\]
Choose the source fixed bump $\widehat\rho$ in its reserved mean
interval $I_m$, with $\int\widehat\rho(x)\,dx=1$, and put
\[
 \rho_{\rm phys}=q^{-1/2}\widehat\rho(r/\sqrt q),\qquad
 c_{\rho,\rm phys}=\frac12\int r^2\rho_{\rm phys}\,dr
             =\frac q2\int x^2\widehat\rho(x)\,dx .
\]
Thus $(c_{\rho,\rm phys})_z=(q_z/q)c_{\rho,\rm phys}$;
it is generally not constant in $z$.

For precision the source's compact physical primitive is as follows.
The fixed slow cutoff $\chi_m$ is zero on the inner collar and one
on the outer collar of the active annulus. For an auxiliary-dependent
scalar $f$ define
\[
 \begin{split}
 I f(r,Y)&=\int_0^r
 f(s,Y+(s^{d_r}-r^{d_r})v_r)\,ds,\\
 J f(r,Y)&=\int_0^\infty
 f(s,Y+(s^{d_r}-r^{d_r})v_r)\,ds,\qquad
 I_cf=If-\chi_mJf.
 \end{split}
\]
Here $d_r=2((1+h)\rho_g-h\kappa_s)$ is the exact source exponent,
$\rho_g=\log(4-\sqrt2)/\log(4+\sqrt2)$,
$v_r=(1,-(\sqrt2-1))$, and $\kappa_s=10^{-5}$.
The derivative $\mathcal D_r$ of the shifted integral differentiates
the integration endpoint while cancelling the derivative of its
$-r^{d_r}v_r$ shift. Consequently
$\mathcal D_rI_cf=f-(\partial_r\chi_m)Jf$.
Also $\overline{Jf}=\int\overline f\,dr$, by translation invariance
of Haar measure. This proves the mean identities used next directly.

The source reconstruction after the actual wave addition changes the
physical mean pressure by
\[
 \Delta p_m=I_c(\Delta g_r-\rho_{\rm phys}\Delta P).
\]
Its source has zero radial auxiliary-mean integral by
\eqref{momop:deltaP}. Therefore
\begin{equation}\label{momop:pressureidentity}
 \partial_r\overline{\Delta p_m}
     =g-\rho_{\rm phys}\Delta P,\qquad
 \int r\,\overline{\Delta p_m}\,dr
     =-\frac12\int r^2g\,dr+c_{\rho,\rm phys}\Delta P .
\end{equation}
The first identity follows from the displayed primitive formula;
its cutoff remainder has zero auxiliary average. The second integrates
$\partial_r(r^2\overline{\Delta p_m})$ between the two zero endpoints.
The full radial residual increment after this pressure update is
$-\rho_{\rm phys}\Delta P-
(\partial_r\chi_m)J(\Delta g_r-\rho_{\rm phys}\Delta P)$.
The second term is retained. The source proves its terminal flatness
under its all-order uniform class bounds; finite smoothness alone is
not that proof.

Use different symbols for the source's two flux defects and the raw
candidate fluxes $J_2,J_1$. Their increments are exactly
\begin{align}
 \Delta J_\theta&=J_2,\nonumber\\
 \Delta J_z&=J_1-\frac12\int r^2g\,dr\nonumber\\
 &=\int r\left(C_{zz}-\frac12C_{rr}
                         -\frac12C_{\theta\theta}\right)\,dr
                    +\frac12\partial_z\int r^2C_{zr}\,dr .
 \label{momop:defectdictionary}
\end{align}
These are the differences of source (8.15): the mean-velocity terms
are unchanged and the full wave covariance changes by $C$.
The axial pressure identity \eqref{momop:pressureidentity} now proves
the exact moments after pressure reconstruction:
\begin{equation}\label{momop:afterpressure}
 M_{2,\rm new}=\partial_z\Delta J_\theta=M_2,\qquad
 M_{1,\rm new}=\partial_z(\Delta J_z+c_{\rho,\rm phys}\Delta P)
  =M_1+\partial_z\left(-\frac12\int r^2g\,dr
                                      +c_{\rho,\rm phys}\Delta P\right).
\end{equation}
These are increment identities. The old source moments and defects
are added to obtain the full updated state. No hypothesis that the
old defects vanish was used: identical old mean-velocity terms
cancel in the subtraction. The product involving $c_{\rho,\rm phys}$
stays inside $\partial_z$, including its nonzero derivative.

\subsection{New quantitative budgets for the three actual inputs}

Retain the previously proved quantities
$\mathcal W_e,\mathcal W_{e,z},\mathcal V_e,\mathcal Z_e$ for
$e=1,2$. They satisfy
\[
 \|w_*\|_{[e]}=\mathcal W_e,\quad
 \|\partial_zw_*\|_{[e]}=\mathcal W_{e,z},\quad
 \|w_y\|_{[e]}\le\mathcal V_e,\quad
 \|\partial_zw_y\|_{[e]}\le\mathcal Z_e.
\]
In particular $\mathcal V_e,\mathcal Z_e$ retain separately the
actual amplitude bounds
$Y_1=(A_0/\mu)\widehat\lambda_3^{-7/8}$ and
$Y_2=2A_0\widehat\lambda_3^{1/8}/(\sigma_2\sqrt{n_i})$
and their physical source-band factors. Define merely for the
following displayed estimates
\[
 B_e=2\mathcal W_e\mathcal V_e+\mathcal V_e^2,\qquad
 D_e=\mathcal W_{e,z}\mathcal V_e+
                    \mathcal W_e\mathcal Z_e+\mathcal V_e\mathcal Z_e.
\]
Choose $r_->0$ enclosing the radial support of the new fields and
their axial derivatives from below. Cauchy--Schwarz on the original
weighted product measures gives
\begin{equation}\label{momop:targetbounds}
 |\Delta P|\le r_-^{-2}B_1+2r_-^{-1}D_1,\qquad
 |\Delta J_\theta|\le B_2,\qquad
 |\Delta J_z|\le B_1+D_2.
\end{equation}
Here is the full estimate. The matrix
$\operatorname{diag}(-1,1,0)$ has operator norm one, so
$|C_{\theta\theta}-C_{rr}|$ is bounded after averaging by
$2|w_*||w_y|+|w_y|^2$. Replacing $r^{-1}$ by
$r_-^{-2}r$ and applying Cauchy--Schwarz gives the first part
$r_-^{-2}B_1$ of the pressure bound. Differentiating $C_{zr}$
produces the six actual cross and self product terms. Their integral
with weight $r^e$ is at most $2D_e$. Replace the unweighted
integral by $r_-^{-1}$ times the $r$-weighted integral to obtain
$2r_-^{-1}D_1$ in \eqref{momop:deltaP}.
The same product estimate without differentiation bounds
$\int r^2C_{z\theta}$ by $B_2$.
Finally $\operatorname{diag}(-1/2,-1/2,1)$ has norm one.
Apply it to the first integral in \eqref{momop:defectdictionary}
to get $B_1$; the last differentiated integral is bounded by
$\tfrac12(2D_2)=D_2$. This proves all three claims without
replacing the two amplitude components by a new unspecified constant.

\subsection{The exact five-row map for these targets}

The source map acts on the \emph{current} physical triple
$(P,J_\theta,J_z)$, after adding \eqref{momop:deltaP} and
\eqref{momop:defectdictionary} to the old source triple and
reconstructing pressure. It does not take
$(M_2,M_1)$ as its input. The five equations, in physical variables, are
\begin{align}
 \int r^2\Delta v\,dr&=0,&
 \int r\gamma_d\,dr&=0,\nonumber\\
 \int\frac{2V}{r}\Delta v\,dr&=-P,\nonumber\\
 \int r^2(G\Delta v+V\gamma_d)\,dr&=-J_\theta,\nonumber\\
 \int(2rG\gamma_d-rV\Delta v)\,dr&=-J_z .
 \label{momop:fivephysical}
\end{align}
In the last row there is exactly one factor $r$ multiplying $V\Delta v$.
The pressure-source term contributes
$-\tfrac12\int r^2(2V/r)\Delta v\,dr=-\int rV\Delta v\,dr$,
which proves that factor rather than treating the fifth row as an
independent constraint with an arbitrary weight.

Retain the source's reserved patch $x=r/\sqrt q\in I_m$. Its actual
summed base satisfies source (5.44) and (8.24):
\[
\begin{gathered}
 V_{\rm phys}=q^{-A}a(\eta)x^{-1-2\lambda},\qquad
 G_{\rm phys}=0,\\
 a(\eta)=2^{1/2+\lambda}c_{\rm patch}(1+\eta^2)^{-1},
 \quad \lambda>0,\quad |a(\eta)|\ge a_0>0.
\end{gathered}
\]
The source fixed $\lambda$ here is its reserved-profile exponent,
not a frequency of the coupled problem. In a fixed $Q$ chart,
the same physical base is
\[
 V_*=(Q/q)^A a(\eta)\bigl(\sqrt{Q/q}\,R\bigr)^{-1-2\lambda},
 \qquad G_*=0.
\]
The factor $(Q/q)^A$ has been retained.

Choose a nonnegative nonzero bump $\varphi_0\in C_c^\infty(I_m)$
supported near an interior point and a fixed spacing $d_b>0$ so small
that the three profiles
\[
 \varphi_j(x)=e^{-jd_b}\varphi_0(e^{-jd_b}x),\qquad j=0,1,2,
\]
have disjoint supports in the patch. This is possible because the patch
contains an open interval about that point: choose three distinct
scaled centers there, then decrease the support radius below half the
minimum separation. The symbol $d_b$ distinguishes this source bump
spacing from the coupled physical diffusivity $d$.
For every real $p$, set
\[
 \mu_p=\int x^p\varphi_0(x)\,dx>0,\qquad
 \int x^p\varphi_j(x)\,dx=\mu_pe^{jd_bp}.
\]
The equality is the substitution $x=e^{jd_b}\xi$ and retains the
prefactor in $\varphi_j$.

Define the two explicit moment matrices
\[
 A_\theta=
 \begin{pmatrix}
 \mu_2&\mu_2e^{2d_b}&\mu_2e^{4d_b}\\
 2a\mu_{-2-2\lambda}&
 2a\mu_{-2-2\lambda}e^{d_b(-2-2\lambda)}&
 2a\mu_{-2-2\lambda}e^{2d_b(-2-2\lambda)}\\
 -a\mu_{-2\lambda}&-a\mu_{-2\lambda}e^{-2d_b\lambda}&
                              -a\mu_{-2\lambda}e^{-4d_b\lambda}
 \end{pmatrix},
\]
\[
 A_z=
 \begin{pmatrix}
 \mu_1&\mu_1e^{d_b}\\
 a\mu_{1-2\lambda}&a\mu_{1-2\lambda}e^{d_b(1-2\lambda)}
 \end{pmatrix},\qquad a=a(\eta).
\]
Both inverses exist at every source parameter value. Indeed the
angular matrix is a diagonal row factor times the matrix with rows
$(1,t_p,t_p^2)$ for $p=2,-2-2\lambda,-2\lambda$ and
$t_p=e^{d_bp}$. Subtracting the first row from the other two
and factoring their differences gives
\[
 \det[1,t_p,t_p^2]_{p=p_1,p_2,p_3}
 =(t_{p_2}-t_{p_1})(t_{p_3}-t_{p_1})(t_{p_3}-t_{p_2}).
\]
All three powers are distinct for $\lambda>0$. The diagonal
row factor has determinant
\[
-2a^2\mu_2\mu_{-2-2\lambda}\mu_{-2\lambda}\ne0.
\]
Likewise
\[
\det A_z=a\mu_1\mu_{1-2\lambda}
(e^{d_b(1-2\lambda)}-e^{d_b})\ne0.
\]
This also gives an explicit finite inverse norm bound
$\|A^{-1}\|\le\|\operatorname{adj}A\|/|\det A|$
for either matrix. No uniform limit as $\lambda\downarrow0$ is taken.

Use the exact row normalizations
\[
 P_q=q^{2A}P,\qquad
 J_{\theta,q}=q^{2A-3/2}J_\theta,\qquad
 J_{z,q}=q^{2A-1}J_z,
\]
and define the five coefficients by
\begin{equation}\label{momop:coefficients}
 {\bf u}=A_\theta^{-1}(0,-P_q,-J_{z,q})^T,\qquad
 {\bf s}=A_z^{-1}(0,-J_{\theta,q})^T .
\end{equation}
Then the actual physical velocity increments are
\begin{equation}\label{momop:fiveincrement}
 \Delta v=q^{-A}\sum_{j=0}^2u_j\varphi_j(r/\sqrt q),
 \qquad
 \gamma_d=q^{-A}\sum_{j=0}^1s_j\varphi_j(r/\sqrt q).
\end{equation}
The first and second velocity moments have factors $q^{3/2-A}$
and $q^{1-A}$, respectively. The three inhomogeneous rows have
factors $q^{-2A}$, $q^{3/2-2A}$, and $q^{1-2A}$, respectively.
Substituting the original base into \eqref{momop:fivephysical}
and changing variables $r=\sqrt q\,x$ therefore gives exactly
the two linear systems \eqref{momop:coefficients}. This proves all
five equations, uniqueness in the selected five-dimensional span,
and smooth dependence on the actual target values. For instance
$\partial_\xi A^{-1}=-A^{-1}(\partial_\xi A)A^{-1}$ follows by
differentiating $AA^{-1}=I$; repeated product rules give every
coefficient derivative with the actual $q,\eta$ factors retained.

The desired axial increment has zero weighted integral and is
auxiliary independent. Its exact compact azimuthal potential and
radial increment are
\[
 \Psi(r,z,t)=\frac1r\int_0^r s\gamma_d(s,z,t)\,ds,\qquad
 \Delta\beta=-\partial_z\Psi,\qquad
 \Delta\gamma=(\partial_r+r^{-1})\Psi=\gamma_d.
\]
The zero moment in \eqref{momop:fivephysical} makes $\Psi$ vanish
beyond its outer support; its lower-end definition makes it vanish
below the inner support. Differentiation proves
$(\partial_r+r^{-1})\Delta\beta+\partial_z\Delta\gamma=0$.
The support includes the intervals between bump supports, which still
lie within the reserved patch.
To display the axial scaling cost explicitly, put
$\Psi_j(x)=x^{-1}\int_0^x \xi\varphi_j(\xi)\,d\xi$. Then
\begin{align}
 \Psi&=q^{1/2-A}\sum_{j=0}^1s_j\Psi_j(x),\nonumber\\
 \Delta\beta
 &=-q^{1/2-A}\sum_{j=0}^1\left[
    (\partial_zs_j)\Psi_j+
    \frac{q_z}{q}s_j\left((1/2-A)\Psi_j
                                      -\frac x2\Psi_j'\right)\right].
 \label{momop:inducedradial}
\end{align}
The two individual $\Psi_j$ can have $x^{-1}$ tails. The exact
linear relation $\sum_js_j\mu_1e^{jd_b}=0$ cancels those tails
in their sum and, after differentiation, in \eqref{momop:inducedradial}.
No individual compactness was assumed.

\subsection{What the five-row operation actually leaves}

Let $(\beta,v,\gamma)$ be the current mean correction, after whatever
preceding wave and mean updates have actually been performed. Keep the
base $(b,V,G)$ and the full current wave $w$ fixed during this
five-row step. Use the three full current defects as its targets.
Write $\mathcal D_r$ for the physical extended radial derivative and
$\mathcal D_t$ for the physical extended time derivative. With
$\Delta_0=\mathcal D_r^2+r^{-1}\mathcal D_r+\partial_z^2$, direct
subtraction of the full radial source gives
\begin{align}
 R_g:={}&g_{r,\rm new}-g_r-\frac{2V}{r}\Delta v\nonumber\\
 ={}&-\mathcal D_t\Delta\beta
 -(\mathcal D_r+r^{-1})
             \bigl(2b\Delta\beta+2\beta\Delta\beta+(\Delta\beta)^2\bigr)
 \nonumber\\
 &-\partial_z\bigl(b\Delta\gamma+G\Delta\beta+
           \beta\Delta\gamma+\gamma\Delta\beta+\Delta\beta\Delta\gamma\bigr)
 \nonumber\\
 &+\frac{2v\Delta v+(\Delta v)^2}{r}
               +\nu_{\rm NS}(\Delta_0-r^{-2})\Delta\beta .
 \label{momop:Rg}
\end{align}
Here all increments are auxiliary independent, so their
$\mathcal D_t$ is ordinary $\partial_t$. The products with the
old means still retain their auxiliary dependence. Equation
\eqref{momop:Rg} keeps the actual physical viscosity; the source
unit-viscosity formula (8.26) is its case $\nu_{\rm NS}=1$.

The full updated defects, with pressure recomputed from the new
radial source, are exactly
\begin{align}
 P_{\rm new}&=\int\overline{R_g}\,dr,\nonumber\\
 (J_\theta)_{\rm new}
    &=\int r^2\overline{\gamma\Delta v+v\Delta\gamma+
                                      \Delta\gamma\Delta v}\,dr,\nonumber\\
 (J_z)_{\rm new}
    &=\int r\overline{2\gamma\Delta\gamma+(\Delta\gamma)^2}\,dr
                      -\frac12\int r^2\overline{R_g}\,dr .
 \label{momop:recomputed}
\end{align}
To verify these formulas, expand every quadratic product in the
definitions of the current and new defects. The radial source
increment is $(2V/r)\Delta v+R_g$. The third row of
\eqref{momop:fivephysical} cancels the old $P$ with the integral
of the displayed linear term. The azimuthal flux increment contains
the base-linear terms $G\Delta v+V\Delta\gamma$, cancelled by
the fourth row because $\Delta\gamma=\gamma_d$. The axial defect
contains $2G\Delta\gamma$ and the radial-source contribution
$-\tfrac12r(2V\Delta v)$ after integration. These are precisely
the fifth row. All remaining products are the terms displayed
in \eqref{momop:recomputed}; the full wave covariance is unchanged
during this mean step. This proves the formulas without assuming
their right-hand sides vanish.

The two velocity moments remain exactly zero. The new residual
moments therefore have the source form
\[
 \int r^2\overline{E_{\theta,\rm new}}\,dr
       =\partial_z(J_\theta)_{\rm new},\qquad
 \int r\overline{E_{z,\rm new}}\,dr
       =\partial_z\bigl((J_z)_{\rm new}
                           +c_{\rho,\rm phys}P_{\rm new}\bigr).
\]
The pressure and $c_{\rho,\rm phys}$ are the actual reconstructed
objects. Thus the five equations cancel specified linear defect
contributions and expose the remaining nonlinear and differential
terms. Equations \eqref{momop:Rg}--\eqref{momop:recomputed} supply
the next actual inputs rather than a claim of complete cancellation.

\subsection{Exact source operation dictionary and the bounds it requires}

For the source's estimates, $Q=2^{-\ell}$, $\varepsilon=Q^h$,
$S_*=\ell^2$, $R=r/\sqrt Q$, $Z=z/Q^D$, $T=(1-t)/Q$.
Use a coarsest applicable common auxiliary torus $H$ as in source
Lemma 6.2; different representatives agree by its integer covering
maps. The original radial edge weights are
\[
\begin{gathered}
 \zeta(X)=\exp\left(-\frac{a_a}{\log^2(X/X_a)}
                       -\frac{a_b}{\log^2(X_b/X)}\right),\\
 \delta(X)=\min\{1,\log(X/X_a),\log(X_b/X)\},
 \quad X=\frac{QR^2}{2q},
\end{gathered}
\]
with the source fixed $a_a,a_b>0$. In the following bounds the band,
label, representative and harmonic are held fixed when differentiating.
The constants may depend on the finite correction stage and derivative
multi-index, but are uniform over bands, labels, points and rectangle
copies.

A mean coefficient $f\in M^\alpha$ has no angular dependence, descends
to the common torus, extends smoothly by zero outside the active shell,
and satisfies for each multi-index $I$ in $(R,Z,T,H)$
\[
 |\partial^If|\le C_I\varepsilon^\alpha S_*^{b_I}
                                        \zeta\,\delta^{-d_I}.
\]
It can have nonzero auxiliary modes and has no rectangle or pulse
envelope restriction. A moment coefficient $F(Z,T)\in S^\alpha$
satisfies $|\partial_{Z,T}^IF|\le C_I\varepsilon^\alpha S_*^{b_I}$.
A labelled nonzero harmonic amplitude $a_{\gamma,m}\in W^\alpha$
has no angular dependence, the source slow, transverse and shell
support, smooth zero extension across those boundaries, compatible
lift representatives, and
\[
 |\partial^Ia_{\gamma,m}|\le
 C_I\varepsilon^\alpha S_*^{b_I}\sqrt\zeta\,\delta^{-d_I}
                          {\cal P}_\gamma(v),\qquad 0\le v\le L_s.
\]
The prescribed source phase has $k_\gamma p_\gamma\in
\mathbb Z\setminus\{0\}$ and $k_\gamma=\lceil\varepsilon^{-1/2}\rceil$.
The harmonic set is finite at a fixed stage and independent of band.
Equal label--harmonic terms are grouped before testing the bounds.
The pulse envelope is the original positive exponential in (7.12),
with its proved Gaussian estimates. Actual cut-off waves additionally
have support where the source pulse cutoff $\psi$ is supported and
smooth zero extension in $v$. No divisibility by $\psi$ is required.

The source's anisotropic derivative mappings, with the phase factored
out, are
\[
\begin{gathered}
 \partial_{R,Z,T}:C^\alpha\to C^\alpha,\quad
 D_r:C^\alpha\to C^{\alpha-\kappa_s},\quad
 D_z=\varepsilon\partial_Z:C^\alpha\to C^{\alpha+1},\\
 c_{i_0}N_{i_0}:C^\alpha\to C^\alpha,\quad
 -\varepsilon\partial_T:C^\alpha\to C^{\alpha+1},
 \qquad C=M,W.
\end{gathered}
\]
Products obey $M^\alpha M^\beta\subset M^{\alpha+\beta}$ and
$M^\alpha W^\beta\subset W^{\alpha+\beta}$. Same-label wave
products have this summed order, in $M$ for zero output harmonic
and $W$ otherwise; distinct supported labels have zero products.
The zero-harmonic global $M$ statement here is for the actual temporally
cut-off waves. Before that cutoff, the source gives the mean bounds on
the labelled rectangle, not a global $M$ assertion.
These are source Proposition 6.6, not deductions from the candidate's
finite norm bounds.

Every radial target must first carry the correct measure and weight.
For $f_*=Q^{a_f}f_{\rm phys}$,
\[
 \int R^e\overline{f_*}\,dR
       =Q^{a_f-(e+1)/2}\int r^e\overline{f_{\rm phys}}\,dr .
\]
For the present inputs this reads
\begin{align*}
 P_*&=Q^{2A}P,&
 (J_\theta)_*&=Q^{2A-3/2}J_\theta,&
 (J_z)_*&=Q^{2A-1}J_z,&
 (c_\rho)_*&=Q^{-1}c_{\rho,\rm phys},\\
 (M_2)_*&=Q^{2A-1}M_2,&
 (M_1)_*&=Q^{2A-1/2}M_1 .
\end{align*}
In particular \eqref{momop:afterpressure} becomes exactly
$\varepsilon\partial_Z\Delta J_{\theta,*}$ and
$\varepsilon\partial_Z(\Delta J_{z,*}+c_{\rho,*}\Delta P_*)$.
The factors match because $h+D=1/2$. This is the extra axial
factor used in source (9.12); using the raw $J_1$ after pressure
reconstruction would omit part of its actual input.

For clarity, the following are the four ordered operations of source
Proposition 9.6, with the full current velocity and pressure
reconstructed after each operation.

\begin{enumerate}
\item For each supported nonzero harmonic coefficient $f_{\gamma,m}$
of the full current residual, solve the original transverse pulse
equation with zero entrance data and source $-f_{\gamma,m}$.
In frame coordinates its coefficient matrix is
$H_\gamma-m^2\varepsilon k_\gamma^2|n_\gamma|^2I$.
Multiply the solved potential and pressure by $\psi$ and take the
full curl. Source Proposition 7.2 maps $W^\alpha$ sources to
$W^\alpha$ velocity amplitudes and $W^{\alpha+1/2}$ pressure.
Proposition 9.1 assigns the remaining supported linear residual
order $\alpha+1/2-3\kappa_s$ and retains its Gaussian cutoff tails
additively. The actual candidate's full residual formula determines
the input coefficients. Following this wave change, replace $w_y$
by the actual total new wave increment in \eqref{momop:C}, recompute
$g_r$ and reconstruct pressure. Its old-wave cross terms must change
with the updated wave.

\item Let $F_2,F_1$ be the auxiliary-averaged physical tangential
residuals of that updated state. Their two moments are the updated
versions of \eqref{momop:afterpressure}, including the old defects.
For each $e=1,2$ retain the fixed normalized physical bump
$b_e=q^{-(e+1)/2}\widehat b_e(r/\sqrt q)$ with
$\int x^e\widehat b_e(x)\,dx=1$. Form
\[
 \mathfrak M_e=\int r^eF_e\,dr,\qquad
 \sigma_e=-r^{-e}\int_0^r s^e
                         (F_e(s)-b_e(s)\mathfrak M_e)\,ds .
\]
The zero total weighted integral proves compact support and
$(\partial_r+e/r)\sigma_e=-F_e+b_e\mathfrak M_e$.
Use $\Sigma=Q^{2A}(\sigma_2,\sigma_1)$ as the input to the
source's signed amplitude derivative (7.31), with its fixed
positive primary amplitudes. Source Proposition 7.6 maps an
auxiliary-independent $M^\alpha$ stress to $W^{\alpha-1/2}$.
Take its actual curl, retain the old-wave cross and quadratic
correction tensor (9.13), and reconstruct pressure. The source
does not set $b_e\mathfrak M_e$ to zero.

If $\mathfrak M_e=\partial_z\mathfrak J_e$, its bump term satisfies
the exact derivative identity
\[
 b_e\partial_z\mathfrak J_e
 =\partial_z(b_e\mathfrak J_e)-(\partial_zb_e)\mathfrak J_e,\qquad
 \partial_zb_e=-\frac{q_z}{2q}
                    \bigl((e+1)b_e+r\partial_rb_e\bigr).
\]
It follows by differentiating the original $q$-scaled bump.
Thus trading the moment bump for an axial tensor flux leaves the
displayed radial-profile derivative term. For the axial equation,
$\mathfrak J_1=J_z+c_{\rho,\rm phys}P$, including both product
derivatives in \eqref{momop:afterpressure}.

\item For the zero-auxiliary-mean part of the current angularly
invariant residual, first convert the physical residual to the fixed
source chart. In this step a star denotes that chart representative:
\[
\begin{gathered}
 E_*=Q^{2A+1/2}E_{\rm phys},\qquad
 E_{a,*}^\circ=E_{a,*}-\langle E_{a,*}\rangle_Y,\\
 (\Delta\beta_*,\Delta v_*,\Delta\gamma_*)
       =Q^A(\Delta\beta_{\rm phys},\Delta v_{\rm phys},
                                      \Delta\gamma_{\rm phys}).
\end{gathered}
\]
The source temporal update in these exact units is
\[
 \Delta v_*=-c_{i_0}^{-1}N_{i_0}^{-1}E_{\theta,*}^\circ,\qquad
 \gamma_{d,*}=-c_{i_0}^{-1}N_{i_0}^{-1}E_{z,*}^\circ,
 \quad
 N_{i_0}^{-1}F=
 \sum_{k\ne0}\frac{\widehat F(k)}{2\pi i v_t\cdot k}e^{2\pi ik\cdot H}.
\]
Here $v_t=(\sqrt2-1,1)$ and
$c_{i_0}=T_g^{i_0}Q^{1+h}$, with $c_{i_0}^{-1}\le CS_*$.
The source Diophantine bound gives
$\|N^{-1}F\|_{C^m}\le C_m\|F\|_{C^{m+4}}$.
Let $T_{1,*}$ and $A_{1,*}$ be the source's radial potential and
its cutoff remainder in this chart. Realize $\gamma_{d,*}$ by
$T_{1,*}\gamma_{d,*}$:
$\Delta\beta_*=-\varepsilon\partial_ZT_{1,*}\gamma_{d,*}$ and
$\Delta\gamma_*=\gamma_{d,*}-A_{1,*}\gamma_{d,*}$.
The actual axial remainder $-A_{1,*}\gamma_{d,*}$ stays in the
new velocity. The exact normalized cancellations are
\[
\begin{split}
 c_{i_0}N_{i_0}\Delta v_*&=-E_{\theta,*}^\circ,\\
 c_{i_0}N_{i_0}\Delta\gamma_*&=-E_{z,*}^\circ
              -c_{i_0}N_{i_0}A_{1,*}\gamma_{d,*}.
\end{split}
\]
Convert every resulting velocity component back by $Q^{-A}$ before
reconstructing the physical pressure or taking physical products and
defects. In particular its physical angular update is exactly
\[
 \Delta v_{\rm phys}
  =-Q^{-A}c_{i_0}^{-1}Q^{2A+1/2}
                     N_{i_0}^{-1}E_{\theta,\rm phys}^\circ
  =-T_g^{-i_0}N_{i_0}^{-1}E_{\theta,\rm phys}^\circ .
\]
The last equality uses $A+1/2=1+h$ and the original
$c_{i_0}=T_g^{i_0}Q^{1+h}$. It is also
$-N_{\rm abs}^{-1}E_{\theta,\rm phys}^\circ$, since
$N_{\rm abs}=T_g^{i_0}N_{i_0}$ on the representative.
This proves the conversion without applying a normalized inverse
directly to a physical residual.

\item Apply \eqref{momop:coefficients} to the three defects of this
updated state, use \eqref{momop:fiveincrement} and its actual
radial component \eqref{momop:inducedradial}, and reconstruct
pressure again. The exact remaining defects are
\eqref{momop:Rg}--\eqref{momop:recomputed}.
Source Lemma 8.7 maps $S^\alpha$ targets to $M^\alpha$
tangential increments and $M^{\alpha+1}$ radial increments.
Source Lemma 8.8 gives the remaining defect class
$S^{\alpha+0.9-2\kappa_s}$ using the actual bounds
$v,\gamma\in M^{0.9}$, $\beta\in M^{1.9}$, $\alpha\ge0.9$,
and its reserved-patch base derivative bounds.
\end{enumerate}

The source cycle assumes its full finite correction-state bounds
$w\in W^{1/2}$, $w-w_0^{\tan}\in W^{0.68}$,
$v,\gamma,p_m\in M^{0.9}$, $\beta\in M^{1.9}$,
with supported nonzero residual order
$B_j=1/2+\sigma_j$, mean and defect order $C_j^*=1+\sigma_j$,
and $\sigma_j=1/5+j/10$. Its ordered estimate proves gains
$B_{j+1}-B_j=C_{j+1}^*-C_j^*=1/10$.
These are requirements and conclusions for the source's construction.
The explicit candidate budgets and equations above do not by
themselves identify those exponents.

There is a precise finite-domain statement available now. Write the
full current residual as $R_{\rm old}+\Delta R$, retaining the old
source state and every term of the computed increment. The added
candidate uses finitely many labels, with its cutoffs supported
strictly inside the source rectangle and radial annulus. On those
compact coefficient supports, $\varepsilon>0$, $\zeta>0$,
$\delta>0$ and $\mathcal P_\gamma(v)>0$ have positive minima
for each fixed label. By local finiteness of the retained source-label
family, only finitely many old labels intersect this fixed compact
perturbation support. Every differentiated coefficient of $\Delta R$
is finite by its explicit smooth shapes, including the retained old-state
cross terms there. For any fixed
$\alpha$ and derivative $I$, the maximum of its derivative divided by
$\varepsilon^\alpha\sqrt\zeta\mathcal P_\gamma$ on these finitely
many supports is therefore a concrete finite constant. Extend the
finite \emph{increment} family by zero on unused labels. This proves its
finite coefficient inequalities without assigning an asymptotic gain
to the amplitude. It does not set any old source coefficient to zero:
$R_{\rm old}$ retains its actual locally finite family and its attributed
source-class bounds. The finite-sum closure in source Proposition 6.6
combines those bounds with the proved increment bounds for each full
current input. The same compact-support argument applies to the added
mean sources and the reserved bump operators. Thus the operations
above are defined on the actual old-plus-increment inputs.
Their constants can contain the inverse of a very small Gaussian
envelope or source scale; no uniform control for an increasing family
of coupled stages or bands follows from this compactness argument.
The original component budgets and the full nonlinear remainders
remain the quantities to propagate when investigating such a family.
